\chapter{Rashes in Babies and Children}

\section*{Definition}
A rash is a change in skin color, texture, or appearance. Childhood rashes are often infectious or inflammatory, but some rapidly spreading or non-blanching rashes indicate urgent illness.

\section*{When to See a Doctor}
Call emergency services for breathing difficulty or facial swelling. Obtain urgent care for a purple or non-blanching rash, fever with a child who appears ill, blistering or peeling skin, severe pain, or rapidly spreading rash. Arrange prompt evaluation for persistent, recurrent, unexplained, or function-limiting symptoms.

\section*{How It Develops}
A childhood rash reflects changes in skin blood flow, inflammation, infection, allergy, or irritation. Age, fever, appearance, blanching, mucosal involvement, and the child's general condition help distinguish benign eruptions from emergencies.

\section*{Causes}
\begin{itemize}
  \item Viral exanthems, eczema, hives, contact dermatitis, insect bites, fungal infection, scabies, impetigo, medication reactions, and meningococcal disease.
  \item Medication effects, environmental exposures, and underlying chronic disease may also contribute.
  \item Serious causes are less common but must be considered when symptoms are sudden, progressive, severe, or associated with systemic illness.
\end{itemize}

\section*{Related Symptoms}
\begin{itemize}
  \item Pain, fever, fatigue, nausea, weakness, or changes in normal function
  \item Bleeding, swelling, discharge, breathing difficulty, or neurologic symptoms when a serious cause is present
\end{itemize}

\section*{Medical Evaluation}
\begin{itemize}
  \item History addresses onset, duration, severity, triggers, injuries, exposures, medications, medical conditions, age, and pregnancy status when relevant.
  \item Examination includes vital signs and focused assessment of the relevant organ system, neurologic status, hydration, and function.
  \item Testing may include blood or urine studies, cultures, imaging, physiologic testing, fetal or pediatric assessment, or specialist examination.
\end{itemize}

\section*{Treatment Options}
Treatment depends on the diagnosis: supportive care may be enough for many viral rashes, while eczema, hives, fungal infection, scabies, or bacterial disease needs age-appropriate, cause-specific treatment. Do not apply prescription steroid or antibiotic creams without advice.

\section*{Self-Care and Home Management}
Avoid known triggers, protect the affected area, maintain appropriate hydration, and record timing and associated symptoms. Use only age- or pregnancy-appropriate medicines recommended by a clinician. Do not delay emergency care while trying home treatment.

\section*{References}
\begin{thebibliography}{99}
\bibitem{rashes_in_children:nice} National Institute for Health and Care Excellence. Fever in under 5s: assessment and initial management. NICE guideline NG143; 2019 (updated 2021).
\bibitem{rashes_in_children:paller} Paller AS, Mancini AJ. \textit{Hurwitz Clinical Pediatric Dermatology: A Textbook of Skin Disorders of Childhood and Adolescence}. 6th ed. Elsevier; 2021.
\end{thebibliography}
